\scanpage{248}
\S\ 1 Un théorème de géométrie analytique.

Rappelons d'abord deux résultats connus

\uline{Lemme} 1.1. Soit $X$ un espace topologique, $Y$ une partie fermée de $X$ ayant un syst. fondamental de voisinages paracompacts, $F$ un faisceau abélien sur $X$, $F|Y$ le faisceau induit sur $Y$. Alors l'hom. can. \deleted{est} suivant est un isom :
\[
\varinjlim_U H^*(U,F)\longrightarrow H^*(Y,F|Y)
\]
[où dans le premier mb, le \uline{lim} est \unsure{sur|pris sur} l'ensemble des voisinages de $Y$ de $X$]

\uline{Corollaire} 1.2 Soit $f:X\longrightarrow Y$ une application continue \deleted{propre}\inserted{fermée}, avec $Y$ loc. compact, \deleted{$F$ un} $y\in Y$, $X_0=f^{-1}(\{y\})$, $F$ un faisceau abélien sur $X$, $F|X_0$ le faisceau induit sur $X_0$. Alors on a un isom canonique
\[
R^*\unsure{f(F)|f_*(F)}_y\simeq H^*(X_0,F|X_y)
\]
\mspage{247}
